% TOP_PROBLEM: {"problemId": "problem.sums-of-three-cubes-representation-conjecture", "canonicalTitle": "Sums of Three Cubes Representation Conjecture", "releaseRank": 430, "primaryDomain": "number_theory_arithmetic_geometry", "primaryDomainLabel": "Number theory & arithmetic geometry", "displayStatus": "open_with_solved_subcases", "releaseStatus": "open_with_solved_subcases"}
% !TeX program = lualatex
% Definition and sources only; this file is not an LLM solution attempt.
\documentclass[11pt]{article}
\usepackage[margin=25mm]{geometry}
\usepackage{fontspec}
\setmainfont{Latin Modern Roman}
\usepackage{amsmath,amssymb,mathtools}
\usepackage{unicode-math}
\setmathfont{Latin Modern Math}
\usepackage{xurl,hyperref}
\hypersetup{colorlinks=true,urlcolor=blue}
\setcounter{secnumdepth}{0}
\setlength{\parindent}{0pt}
\setlength{\parskip}{5pt}
\setlength{\emergencystretch}{3em}
\begin{document}
\section*{\texorpdfstring{Sums of Three Cubes Representation Conjecture}{Sums of Three Cubes Representation Conjecture}}\label{430}
\subsection{Definitions and mathematical statement}
For an integer $k$ set
\[
 R_k=\{(x,y,z)\in{\mathbb{Z}}^3:x^3+y^3+z^3=k\}.
\]
Since a cube is congruent to $0,1$, or $-1$ modulo $9$, the congruences $k\equiv4,5\pmod9$ preclude a representation. The basic conjecture is the converse:
\[
 k\not\equiv4,5\pmod9\quad\Longrightarrow\quad R_k\ne\varnothing.
\]
The catalogue also mentions the stronger refinement that $R_k$ is infinite for each eligible $k$; this is a separate strength of assertion. The variables may be negative or zero, and there is no requirement of positivity or pairwise coprimality. A solution whose coordinates are large relative to $k$ is a valid representation.
\par\smallskip\textit{Formulation and concept references:} \href{https://arxiv.org/pdf/2402.07146}{[S1]}.

\subsection{Short English statement}
Is the obstruction modulo nine the only reason an integer cannot be written as three integer cubes? A stronger version asks for infinitely many representations.
\subsection{Sources}
\begin{itemize}
\item[S1]T. Browning, J. Glas, V. Y. Wang, "Sums of three cubes over a function field," arXiv:2402.07146 (2024); published as "Optimal sums of three cubes in F\_q[t]," Math. Z. (2025).
\url{https://arxiv.org/pdf/2402.07146}.
\textit{Formulation source.}
\end{itemize}
\end{document}
